\documentclass[a4paper,11pt]{article}

% ============================================================
%  Le Collé — Séance 1 : Déterminants, courbes paramétrées
%  © 2026 Ahmed Amine Ouqlani
% ============================================================

% ---------- Polices ----------
\usepackage[utf8]{inputenc}
\usepackage[T1]{fontenc}
\usepackage[sc]{mathpazo}          % texte et mathématiques : Palatino
\usepackage[scaled=0.88]{helvet}   % titres et repères : linéale
\linespread{1.07}
\usepackage{amsmath,amssymb,amsthm,mathrsfs,stmaryrd}
\usepackage{microtype}

% ---------- Mise en page ----------
\usepackage[a4paper,top=2.3cm,bottom=2.3cm,left=2.4cm,right=2.4cm,
            headheight=14pt,headsep=0.8cm,footskip=1.1cm]{geometry}
\setlength{\parindent}{0pt}
\setlength{\parskip}{0.45em plus 0.15em minus 0.1em}
\raggedbottom
\emergencystretch=1.5em
\clubpenalty=10000 \widowpenalty=10000 \displaywidowpenalty=10000 \brokenpenalty=10000

% ---------- Couleurs : un bleu de nuit, un ocre, des gris ----------
\usepackage{xcolor}
\definecolor{Ink}{HTML}{1C2533}
\definecolor{Bleu}{HTML}{1F3A5F}
\definecolor{BleuMoyen}{HTML}{4A6A8F}
\definecolor{BleuClair}{HTML}{9DB2C9}
\definecolor{Fond}{HTML}{F2F4F7}
\definecolor{Ocre}{HTML}{B07A26}
\definecolor{Brique}{HTML}{A23B2A}
\definecolor{Gris}{HTML}{5E6672}
\definecolor{Filet}{HTML}{C5CCD6}

% ---------- Packages ----------
\usepackage{enumitem}
\setlist{itemsep=0.25em,topsep=0.35em,parsep=0pt}
\setlist[itemize]{label=\textcolor{BleuMoyen}{\small\textbullet},leftmargin=1.4em}
\setlist[enumerate]{leftmargin=1.8em,label=\textsf{\small\color{Bleu}\arabic*.}}
\usepackage{booktabs,array,multicol}
\usepackage{needspace,float}
\usepackage[font=small,labelfont={sf,color=Bleu},labelsep=period,skip=6pt]{caption}
\usepackage{tikz}
\usetikzlibrary{arrows.meta,calc,decorations.markings}
\usepackage{pgfplots}
\pgfplotsset{compat=1.17}
\usepackage{tkz-tab}
\usepackage[most]{tcolorbox}
\usepackage{titlesec,titletoc}
\usepackage{fancyhdr,lastpage,etoolbox}
\usepackage[colorlinks=true,linkcolor=Bleu,urlcolor=Bleu,citecolor=Bleu,
            pdftitle={Le Collé — TD Algèbre linéaire (énoncés)},pdfauthor={Ahmed Amine Ouqlani}]{hyperref}

\numberwithin{equation}{section}
\numberwithin{figure}{section}
\numberwithin{table}{section}

% ---------- Titres ----------
% Chapitres (sections) : grand numéro, titre en linéale, filet ocre
\titleformat{\section}[display]
  {\normalfont\sffamily\color{Ink}}
  {{\small\color{Gris}\textls[200]{\MakeUppercase{Chapitre}}}\\[-0.2em]%
   {\fontsize{64}{64}\selectfont\color{BleuClair}\thesection}}
  {0.4em}
  {\fontsize{25}{30}\selectfont\bfseries}
  [\vspace{0.5em}{\color{Ocre}\rule{3.2cm}{1.4pt}}]
\titlespacing*{\section}{0pt}{0pt}{1.6em}
\titleformat{name=\section,numberless}[display]
  {\normalfont\sffamily\color{Ink}}{}{0pt}{\fontsize{25}{30}\selectfont\bfseries}
  [\vspace{0.5em}{\color{Ocre}\rule{3.2cm}{1.4pt}}]
\titlespacing*{name=\section,numberless}{0pt}{0pt}{1.6em}

\titleformat{\subsection}
  {\needspace{5\baselineskip}\normalfont\sffamily\Large\bfseries\color{Ink}}
  {\color{Bleu}\thesubsection}{0.8em}{}
\titlespacing*{\subsection}{0pt}{1.8em}{0.5em}
\titleformat{\subsubsection}
  {\needspace{5\baselineskip}\normalfont\sffamily\large\bfseries\color{Ink}}
  {\color{Bleu}\thesubsubsection}{0.7em}{}
\titlespacing*{\subsubsection}{0pt}{1.6em}{0.4em}
\titleformat{\paragraph}[runin]{\normalfont\sffamily\bfseries\color{Ink}}{}{0pt}{}[.]
\titlespacing*{\paragraph}{0pt}{0.9em}{0.6em}
\setcounter{secnumdepth}{3}
\setcounter{tocdepth}{2}

% ---------- Table des matières ----------
\titlecontents{section}[0pt]{\addvspace{0.7em}\sffamily\bfseries\color{Ink}}
  {\makebox[2.2em][l]{\color{Bleu}\thecontentslabel}}{}{\hfill\contentspage}[\addvspace{0.2em}]
\titlecontents{subsection}[2.2em]{\small\linespread{1}\selectfont}
  {\makebox[2.6em][l]{\color{Gris}\thecontentslabel}}{}{\titlerule*[0.7pc]{.}\contentspage}

% ---------- En-têtes et pieds de page ----------
\pagestyle{fancy}
\fancyhf{}
\renewcommand{\sectionmark}[1]{\markboth{\thesection\quad#1}{}}
\fancyhead[L]{\footnotesize\sffamily\color{Gris}\leftmark}
\fancyhead[R]{\footnotesize\sffamily\scshape\color{Gris}Le Collé}
\fancyfoot[C]{\small\sffamily\color{Ink}\thepage}
\fancyfoot[L]{\scriptsize\sffamily\hyperlink{toc}{\color{Filet}sommaire}}
\renewcommand{\headrulewidth}{0.4pt}
\renewcommand{\headrule}{{\color{Filet}\hrule height\headrulewidth width\headwidth}}
\renewcommand{\footrulewidth}{0pt}
\fancypagestyle{plain}{\fancyhf{}\fancyfoot[C]{\small\sffamily\color{Ink}\thepage}\renewcommand{\headrulewidth}{0pt}}

% ---------- Énoncés ----------
\newtheoremstyle{enonce}{0pt}{0pt}{}{}{\sffamily\bfseries\color{Bleu}}{.}{0.5em}
  {\thmname{#1}\thmnumber{~#2}\thmnote{\ \normalfont\sffamily\color{Ink}(#3)}}
\newtheoremstyle{exemple}{0.9em}{0.9em}{}{}{\sffamily\bfseries\color{BleuMoyen}}{.}{0.5em}
  {\thmname{#1}\thmnumber{~#2}\thmnote{\ \normalfont\sffamily\color{Ink}(#3)}}
\newtheoremstyle{remarque}{0.7em}{0.7em}{}{}{\sffamily\itshape}{.}{0.5em}{\thmname{#1}\thmnote{\ (#3)}}
\theoremstyle{enonce}
\newtheorem{definition}{Définition}[section]
\newtheorem{theoreme}[definition]{Théorème}
\newtheorem{proposition}[definition]{Proposition}
\newtheorem{corollaire}[definition]{Corollaire}
\theoremstyle{exemple}
\newtheorem{exemple}[definition]{Exemple}
\theoremstyle{remarque}
\newtheorem*{remarque}{Remarque}

\tcbset{
  defi/.style={enhanced,breakable,colback=white,colframe=white,boxrule=0pt,arc=0pt,
     borderline west={1.2pt}{0pt}{Bleu},left=11pt,right=4pt,top=3pt,bottom=3pt,
     before skip=1.1em,after skip=1.1em},
  thm/.style={enhanced,breakable,colback=Fond,colframe=Fond,boxrule=0pt,arc=0pt,
     borderline west={2.4pt}{0pt}{Bleu},left=11pt,right=8pt,top=6pt,bottom=6pt,
     before skip=1.2em,after skip=1.2em}}
\tcolorboxenvironment{definition}{defi}
\tcolorboxenvironment{theoreme}{thm}
\tcolorboxenvironment{proposition}{thm}
\tcolorboxenvironment{corollaire}{thm}

\renewcommand{\proofname}{Démonstration}
\renewcommand{\qedsymbol}{\textcolor{BleuClair}{$\blacksquare$}}
\makeatletter
\renewenvironment{proof}[1][\proofname]{\par\pushQED{\qed}\normalfont
  \topsep0.5em\relax\trivlist\item[\hskip\labelsep\sffamily\itshape\color{Gris}#1.]\ignorespaces}
  {\popQED\endtrivlist\@endpefalse}
\makeatother

% ---------- Méthode : filets haut et bas, titre en bandeau ----------
\makeatletter\newcommand{\methlabelfix}{\protected@edef\@currentlabel{\thetcbcounter}}\makeatother
\newtcolorbox[auto counter,number within=section]{methode}[2][]{%
  enhanced,unbreakable,colback=white,colframe=white,boxrule=0pt,arc=0pt,
  borderline north={0.9pt}{0pt}{Bleu},borderline south={0.9pt}{0pt}{Bleu},
  left=4pt,right=4pt,top=8pt,bottom=8pt,before skip=1.4em,after skip=1.4em,
  before upper={\methlabelfix\phantomsection\addcontentsline{met}{metentry}{\protect\makebox[2.6em][l]{\thetcbcounter}#2}%
     {\sffamily\bfseries\small\color{white}\colorbox{Bleu}{\strut\ MÉTHODE~\thetcbcounter\ }}%
     \ \ {\sffamily\bfseries\color{Bleu}#2}\par\smallskip},#1}

% ---------- Repères pédagogiques : un filet de couleur, un mot-clé ----------
\newcommand{\definirbloc}[3]{% nom, couleur, libellé
  \newtcolorbox{#1}[1][]{enhanced,breakable,blanker,
    borderline west={1.4pt}{0pt}{#2},left=10pt,right=2pt,top=1pt,bottom=1pt,
    before skip=1em,after skip=1em,
    before upper={{\sffamily\bfseries\footnotesize\color{#2}\textls[120]{\MakeUppercase{#3}}}\quad},##1}}
\definirbloc{reflexe}{Bleu}{Réflexe}
\definirbloc{astuce}{Ocre}{Astuce}
\definirbloc{attention}{Brique}{Attention}
\definirbloc{enpratique}{BleuMoyen}{En pratique}
\definirbloc{plusloin}{Gris}{Pour aller plus loin}
\definirbloc{strategie}{Bleu}{Stratégie}
\definirbloc{lecture}{BleuMoyen}{Lire le schéma}
% compatibilité avec les parties reprises
\newenvironment{garde}{\begin{attention}}{\end{attention}}
\newenvironment{complementHP}{\begin{plusloin}}{\end{plusloin}}
\newenvironment{interpretation}{\begin{enpratique}}{\end{enpratique}}

% ---------- Exercices, corrigés ----------
\newcounter{exo}[section]
\renewcommand{\theexo}{\thesection.\arabic{exo}}
\newcommand{\niveau}[1]{{\color{Ocre}\footnotesize
  \ifcase#1\or$\bullet\circ\circ\circ\circ$\or$\bullet\bullet\circ\circ\circ$\or$\bullet\bullet\bullet\circ\circ$%
  \or$\bullet\bullet\bullet\bullet\circ$\or$\bullet\bullet\bullet\bullet\bullet$\fi}}
\newcommand{\serie}[1]{\par\needspace{8\baselineskip}\vspace{1.5em}%
  \noindent{\sffamily\bfseries\color{Bleu}#1}\par\nobreak\vspace{-0.4em}%
  \noindent{\color{Filet}\rule{\linewidth}{0.4pt}}\par\nobreak\vspace{0.1em}}
\newenvironment{exo}[3]{%
  \par\needspace{5\baselineskip}\vspace{0.75em}\refstepcounter{exo}\label{e:#1}%
  \def\exolab{#1}%
  \noindent{\sffamily\bfseries\color{Bleu}Exercice~\theexo}\quad\niveau{#2}%
  \hfill{\footnotesize\sffamily\color{Gris}#3}\par\nobreak\vspace{0.15em}}%
  {\par\nopagebreak\noindent\hfill{\scriptsize\sffamily\hyperref[c:\exolab]{\color{Gris}corrigé p.~\pageref{c:\exolab}}}\par}
\newenvironment{corrige}[1]{%
  \par\needspace{6\baselineskip}\vspace{1em}\phantomsection\label{c:#1}%
  \noindent{\sffamily\bfseries\color{Bleu}Exercice~\ref{e:#1}}\quad
  {\scriptsize\sffamily\hyperref[e:#1]{\color{Gris}énoncé p.~\pageref{e:#1}}}\par\nobreak\vspace{0.2em}}%
  {\hfill\textcolor{BleuClair}{$\square$}\par}
\newcommand{\etape}[1]{\par\smallskip\noindent{\sffamily\bfseries\small\color{BleuMoyen}#1.}\ }
\newenvironment{solution}{\par\noindent{\sffamily\bfseries\small\color{BleuMoyen}Solution.}\ }{\hfill\textcolor{BleuClair}{$\square$}\par}
\newcommand{\chercher}{\par\smallskip\noindent{\color{Filet}\rule[0.5ex]{1.2cm}{0.4pt}}\
  {\footnotesize\sffamily\itshape\color{Gris}laisser chercher l'étudiant avant de poursuivre}\par\smallskip}

% ---------- Guide de séance ----------

\newcommand{\marge}[1]{\leavevmode\marginpar{\raggedleft\scriptsize\sffamily\color{Gris}\textls[100]{\MakeUppercase{#1}}}\ignorespaces}
\newcommand{\Etape}[4]{%
  \subsection[Étape #1 --- #3 {\color{Gris}(#2)}]{#3}\label{etape:#1}%
  \noindent{\small\sffamily\color{Gris}{\color{Ocre}\bfseries Étape #1}\quad·\quad #2\quad·\quad #4}\par\medskip}
\newcommand{\question}[1]{\par\smallskip\noindent{\sffamily\bfseries\small\color{BleuMoyen}Question.}\ \textit{#1}\par\smallskip}
\newcommand{\rappel}[1]{\par\smallskip\noindent\phantomsection\addcontentsline{rap}{rapentry}{#1}%
  {\sffamily\bfseries\scriptsize\colorbox{Fond}{\color{Gris}\strut RAPPEL}}\ {\sffamily\small\color{Gris}#1.}\ }
\newcommand{\pourquoi}{\par\noindent{\sffamily\bfseries\small\color{Ocre}Pourquoi ?}\ }

% Ouverture de chapitre : \ouverture{objectifs}{prérequis}{durée}
\newcommand{\ouverture}[3]{%
  \begin{tcolorbox}[enhanced,colback=Fond,colframe=Fond,arc=0pt,boxrule=0pt,left=12pt,right=12pt,top=10pt,bottom=10pt,
     before skip=1.2em,after skip=1.6em]
  \begin{minipage}[t]{0.56\linewidth}
  {\sffamily\bfseries\small\color{Bleu}OBJECTIFS}\par\vspace{0.3em}\small #1
  \end{minipage}\hfill
  \begin{minipage}[t]{0.38\linewidth}
  {\sffamily\bfseries\small\color{Bleu}PRÉREQUIS}\par\vspace{0.3em}\small #2
  \ifx\relax#3\relax\else\par\vspace{0.6em}{\sffamily\bfseries\small\color{Bleu}DURÉE}\par\vspace{0.2em}\small #3\fi
  \end{minipage}
  \end{tcolorbox}}
% Fin de chapitre
\newenvironment{essentiel}{\par\needspace{12\baselineskip}\vspace{1.5em}%
  \begin{tcolorbox}[enhanced,unbreakable,colback=white,colframe=white,boxrule=0pt,arc=0pt,
   borderline north={1.4pt}{0pt}{Ocre},left=2pt,right=2pt,top=8pt,bottom=4pt]
  {\sffamily\bfseries\color{Ocre}L'ESSENTIEL DU CHAPITRE}\par\medskip}{\end{tcolorbox}}

% ---------- Notations ----------
\newcommand{\K}{\mathbb{K}}
\newcommand{\R}{\mathbb{R}}
\newcommand{\C}{\mathbb{C}}
\newcommand{\N}{\mathbb{N}}
\newcommand{\Z}{\mathbb{Z}}
\newcommand{\Mn}[1]{\mathcal{M}_{n}(#1)}
\newcommand{\GLn}[1]{\mathrm{GL}_n(#1)}
\newcommand{\LE}{\mathcal{L}(E)}
\newcommand{\rg}{\operatorname{rg}}
\newcommand{\vect}{\operatorname{Vect}}
\newcommand{\Mat}{\operatorname{Mat}}
\newcommand{\Id}{\operatorname{Id}}
\newcommand{\sh}{\operatorname{sh}}
\newcommand{\ch}{\operatorname{ch}}
\renewcommand{\th}{\operatorname{th}}
\DeclareMathOperator{\argsh}{argsh}
\newcommand{\vv}[1]{\overrightarrow{#1}}
\newcommand{\ve}[1]{\vec{#1}}
\newcommand{\ient}[2]{\llbracket #1,#2\rrbracket}
\newcommand{\dt}[1]{\left|\begin{matrix}#1\end{matrix}\right|}
\newcommand{\dis}{\displaystyle}

% ---------- Figures ----------
\pgfplotsset{trig format plots=rad}
\pgfplotsset{courbe/.style={axis lines=middle,axis equal image,
  axis line style={Gris,-{Stealth[length=5pt]}},
  xlabel={$x$},ylabel={$y$},ticklabel style={font=\scriptsize\sffamily,color=Gris},
  every axis x label/.style={at={(ticklabel* cs:1)},anchor=west,font=\small},
  every axis y label/.style={at={(ticklabel* cs:1)},anchor=south,font=\small},
  samples=300,smooth,width=0.6\linewidth}}
\tikzset{arcetude/.style={Bleu,line width=1.25pt},
         reste/.style={BleuClair,line width=1pt},
         tgte/.style={Ocre,line width=0.9pt},
         asy/.style={Gris,dashed,line width=0.6pt},
         pt/.style={circle,fill=Ink,inner sep=1.2pt}}

% ---------- Navigation ----------
\makeatletter
\newcommand*\l@rapentry{\@dottedtocline{1}{0em}{0em}}
\newcommand*\l@metentry{\@dottedtocline{1}{0em}{0em}}
\newcommand{\listemethodes}{\@starttoc{met}}
\newcommand{\listerappels}{\@starttoc{rap}}
\makeatother
% onglets de chapitre sur le bord droit
\usepackage{eso-pic}
\newif\ifonglets \ongletsfalse
\newcommand{\dessineronglet}{%
  \edef\ongcol{\ifnum\value{section}=3 Ocre\else Bleu\fi}%
  \edef\ongnum{\the\value{section}}%
  \pgfmathsetmacro{\ongy}{3.2+2.3*(\ongnum-1)}%
  \begin{tikzpicture}[remember picture,overlay]
  \fill[\ongcol] ([xshift=-0.75cm,yshift=-\ongy cm]current page.north east)
     rectangle ([yshift=-\ongy cm-1.9cm]current page.north east);
  \node[text=white,font=\sffamily\bfseries\large] at ([xshift=-0.37cm,yshift=-\ongy cm-0.95cm]current page.north east)
     {\ongnum};
  \end{tikzpicture}}
\AddToShipoutPictureBG{\ifonglets\ifnum\value{section}>0 \dessineronglet\fi\fi}
% blocs de navigation et d'enseignement
\definirbloc{commencer}{Gris}{Si vous commencez ici}
\definirbloc{horsprogramme}{Brique}{Hors programme}
\newcommand{\HP}{\,{\sffamily\bfseries\scriptsize\textcolor{Brique}{[HP]}}}
\newcommand{\source}[1]{{\footnotesize\sffamily\color{Gris}\,(#1)}}
\newcommand{\coupdepouce}[1]{\par\vspace{-0.2em}\noindent{\footnotesize\sffamily\color{Gris}%
  \textbf{Coup de pouce} (à donner si l'étudiant bloque) : #1}\par}
\newcommand{\pointarret}[1]{\par\vspace{1em}\noindent{\color{Filet}\rule[0.5ex]{\linewidth}{0.4pt}}\par\nobreak
  \noindent{\footnotesize\sffamily\color{Ocre}\bfseries POINT D'ARRÊT POSSIBLE}\quad{\footnotesize\sffamily\color{Gris}#1}\par\vspace{0.6em}}
\newcommand{\demarche}[1]{\noindent{\sffamily\bfseries\small\color{BleuMoyen}Démarche.}\ \textit{#1}\par\smallskip}

% ---------- Corrigés : idées, astuces, pièges, à retenir ----------
\newcommand{\choix}[1]{\par\noindent{\sffamily\bfseries\small\color{BleuMoyen}Choix de la méthode.}\ #1\par\smallskip}
\newcommand{\gain}[1]{\par\noindent{\sffamily\bfseries\small\color{Ocre}Astuce.}\ #1\par}
\newcommand{\piege}[1]{\par\noindent{\sffamily\bfseries\small\color{Brique}Piège.}\ #1\par}
\newtcolorbox{retenirbox}{enhanced,breakable,colback=Fond,colframe=Fond,arc=0pt,boxrule=0pt,
  borderline west={2pt}{0pt}{Ocre},left=8pt,right=6pt,top=4pt,bottom=4pt,before skip=0.7em,after skip=0.5em}
\newcommand{\retenir}[1]{\begin{retenirbox}{\sffamily\bfseries\footnotesize\color{Ocre}À RETENIR}\quad #1\end{retenirbox}}
\newcommand{\TD}[1]{{\color{Ocre}\bfseries TD\,#1}\ ·\ }
\newcommand{\q}[1]{\par\smallskip\noindent{\sffamily\bfseries\color{Bleu}#1}\ }
\providecommand{\iddots}{\mathinner{\mkern1mu\raise1pt\hbox{.}\mkern2mu\raise4pt\hbox{.}\mkern2mu\raise7pt\hbox{.}\mkern1mu}}

\providecommand{\Ker}{\operatorname{Ker}}\renewcommand{\Im}{\operatorname{Im}}\providecommand{\tr}{\operatorname{tr}}\providecommand{\Sp}{\operatorname{Sp}}

% ================================================================ TD : mise en page propre au livret d'exercices
\fancyhf{}
\fancyhead[L]{\footnotesize\sffamily\color{Gris}TD · Algèbre linéaire · énoncés}
\fancyhead[R]{\footnotesize\sffamily\scshape\color{Gris}Le Collé}
\fancyfoot[C]{\small\sffamily\color{Ink}\thepage\,/\,\pageref{LastPage}}
\fancypagestyle{plain}{\fancyhf{}\fancyfoot[C]{\small\sffamily\color{Ink}\thepage\,/\,\pageref{LastPage}}\renewcommand{\headrulewidth}{0pt}}
\makeatletter
\newcommand{\partie}[3]{% numéro, titre, description
  \par\needspace{10\baselineskip}\vspace{1.4em}\phantomsection\label{partie:#1}\pdfbookmark[0]{Partie #1}{partie.#1}%
  {\sffamily\small\color{Gris}\textls[200]{PARTIE #1}}\par\nobreak\vspace{0.1em}
  {\sffamily\bfseries\LARGE\color{Ink}#2}\par\nobreak\vspace{0.3em}
  {\color{Ocre}\rule{2.6cm}{1.2pt}}\par\nobreak\vspace{0.2em}
  {\small\color{Gris}#3}\par\nobreak\vspace{0.6em}}
\newcommand{\exercice}[7]{% numéro, difficulté, titre, origine, notions, identifiant, page du corrigé
  \par\needspace{8\baselineskip}\vspace{1.1em}\phantomsection
  \def\@currentlabel{#1}\label{e:#6}\pdfbookmark[1]{Exercice #1}{e.#6}%
  \noindent{\sffamily\bfseries\large\color{Bleu}Exercice~#1}\quad\niveau{#2}\hfill{\footnotesize\sffamily\color{Gris}#4}\par\nobreak
  \vspace{0.1em}\noindent{\sffamily\bfseries\color{Ink}#3}\par\nobreak
  \vspace{0.1em}\noindent{\footnotesize\sffamily\color{Gris}Notions : #5.\hfill\mbox{corrigé p.~#7}\par}\nobreak
  \vspace{0.35em}}
\newcommand{\renvoi}[1]{\par\nopagebreak\vspace{0.2em}\noindent\hfill{\scriptsize\sffamily\color{Gris}#1}\par}
\newcommand{\place}{\par\vspace{1.6cm plus 0.6cm minus 0.6cm}}
\newcommand{\correction}[4]{% numéro, titre, identifiant, renvoi
  \par\needspace{7\baselineskip}\vspace{1.3em}\phantomsection
  \def\@currentlabel{#1}\label{e:#3}\pdfbookmark[1]{Exercice #1}{c.#3}%
  {\color{Filet}\hrule height 0.4pt}\vspace{0.5em}
  \noindent{\sffamily\bfseries\large\color{Bleu}Exercice~#1}\quad{\sffamily\bfseries\color{Ink}#2}\hfill{\scriptsize\sffamily\color{Gris}#4}\par\nobreak
  \vspace{0.4em}}
\newcommand{\sommaireligne}[3]{\noindent\hbox to 2.2em{\hfill\sffamily\color{Bleu}#1.}\enspace\begin{minipage}[t]{0.8\linewidth}\raggedright#2\end{minipage}\hfill\pageref{e:#3}\par}
\makeatother
\setlength{\parskip}{0.5em plus 0.15em minus 0.1em}

\begin{document}
\thispagestyle{plain}
{\sffamily\small\color{Gris}\textls[150]{LE COLLÉ · MATHÉMATIQUES PT · CHAPITRE 1}}\par\vspace{0.2em}
{\sffamily\bfseries\fontsize{26}{30}\selectfont\color{Ink}TD --- Algèbre linéaire}\par\vspace{0.3em}
{\sffamily\Large\color{BleuMoyen}Énoncés}\par\vspace{0.4em}
{\color{Ocre}\rule{3.2cm}{1.4pt}}\par\vspace{0.8em}
{\small Ce livret rassemble 19 exercices, classés en quatre parties, du plus direct au plus exigeant.
Il reprend \textbf{tous les exercices du livret} (proposés à la fin de chaque partie du cours), complétés par des
\textbf{exercices complémentaires}, absents du livret. L'origine de chaque exercice est indiquée en haut à droite :
\emph{« Livret, exercice~1.5 »} ou \emph{« Exercice complémentaire »}. La difficulté va de \niveau{1} à \niveau{5}.
Les corrections sont dans un livret séparé, \texttt{TD\_al\_corrections.pdf} ; la page de chaque correction est indiquée en tête de l'énoncé.}

\vspace{0.6em}{\sffamily\bfseries\color{Bleu}Sommaire}\par\vspace{0.2em}{\small\setlength{\parskip}{1.5pt}\setlength{\columnsep}{1.2cm}\begin{multicols}{2}\raggedcolumns
\vspace{0.5em}\noindent\begin{minipage}[t]{0.86\linewidth}\raggedright\leavevmode\sffamily\bfseries\color{Ink}I. Applications directes du cours\end{minipage}\hfill\pageref{partie:I}\par\vspace{0.15em}
\sommaireligne{1}{Un hyperplan de $\mathbb K^n$ et une base}{aln-1}
\sommaireligne{2}{Un supplémentaire dans $\mathbb R_2[X]$}{aln-2}
\sommaireligne{3}{Sommes directes dans $\mathbb R^3$}{red-somme}
\sommaireligne{4}{Hyperplans de $\mathbb R^4$}{red-hyp}
\sommaireligne{5}{Projecteurs associés}{aln-12}
\sommaireligne{6}{Fonctions paires et impaires}{rdn-17}
\vspace{0.5em}\noindent\begin{minipage}[t]{0.86\linewidth}\raggedright\leavevmode\sffamily\bfseries\color{Ink}II. Exercices classiques et méthodes incontournables\end{minipage}\hfill\pageref{partie:II}\par\vspace{0.15em}
\sommaireligne{7}{Projecteur et symétrie dans $\mathbb R^3$}{al-proj}
\sommaireligne{8}{Deux propriétés de la trace}{red-trace}
\sommaireligne{9}{La trace détermine une matrice}{aln-5}
\sommaireligne{10}{Rang d'une matrice à paramètres}{aln-6}
\sommaireligne{11}{Un projecteur donné par sa matrice}{aln-4}
\sommaireligne{12}{Lire des blocs}{red-stable}
\vspace{0.5em}\noindent\begin{minipage}[t]{0.86\linewidth}\raggedright\leavevmode\sffamily\bfseries\color{Ink}III. Exercices de synthèse et approfondissement\end{minipage}\hfill\pageref{partie:III}\par\vspace{0.15em}
\sommaireligne{13}{Deux sous-espaces trop grands se coupent}{aln-3}
\sommaireligne{14}{Matrices de rang au plus $1$}{aln-7}
\sommaireligne{15}{Une matrice de carré nul}{aln-8}
\sommaireligne{16}{L'endomorphisme $M\mapsto AM-MA$}{al-crochet}
\sommaireligne{17}{Matrices nilpotentes de taille 2}{aln-11}
\sommaireligne{18}{$AB-BA=I$ en dimension infinie}{aln-13}
\vspace{0.5em}\noindent\begin{minipage}[t]{0.86\linewidth}\raggedright\leavevmode\sffamily\bfseries\color{Ink}IV. Exercices difficiles et type concours\end{minipage}\hfill\pageref{partie:IV}\par\vspace{0.15em}
\sommaireligne{19}{Endomorphismes de $\mathbb C^2$ de trace nulle}{aln-10}
\end{multicols}}
\clearpage
\partie{I}{Applications directes du cours}{Sommes directes, supplémentaires, hyperplans et leurs équations, projecteurs.}
\exercice{1}{1}{Un hyperplan de $\mathbb K^n$ et une base}{Exercice complémentaire}{hyperplan, forme linéaire, base}{aln-1}{1}
Soit $n\geq2$, $(e_1,\dots,e_n)$ la base canonique de $\K^n$ et $H=\{(x_1,\dots,x_n)\in\K^n,\ x_1+\dots+x_n=0\}$.
\begin{enumerate}
\item Déterminer la dimension de $H$.
\item Pour $1\leq i\leq n-1$, on pose $f_i=e_i-e_n$. Montrer que $(f_1,\dots,f_{n-1})$ est une base de $H$.
\end{enumerate}
\place
\exercice{2}{1}{Un supplémentaire dans $\mathbb R_2[X]$}{Exercice complémentaire}{hyperplan, supplémentaire, polynômes}{aln-2}{1}
Soit $F=\{P\in\R_2[X],\ P(1)=0\}$. Montrer que $\R_2[X]=\vect(1)\oplus F$ et décomposer un polynôme $P$ selon cette somme.
\place
\exercice{3}{1}{Sommes directes dans $\mathbb R^3$}{Livret, exercice 1.1}{somme directe, décomposition, hyperplan}{red-somme}{2}
\begin{enumerate}
\item Dans $\R^3$, on pose $F=\{(x,y,z),\ x+y+z=0\}$ et $G=\vect\big((1,1,1)\big)$. Montrer que $\R^3=F\oplus G$ et décomposer
$v=(x,y,z)$ selon cette somme.
\item Dans $\R^3$, les droites $D_1=\vect(1,0,0)$, $D_2=\vect(0,1,0)$, $D_3=\vect(1,1,0)$ sont-elles en somme directe ?
\end{enumerate}
\place
\exercice{4}{2}{Hyperplans de $\mathbb R^4$}{Livret, exercice 1.3}{hyperplan, intersection, supplémentaire}{red-hyp}{2}
Dans $\R^4$, on considère $H_1$ : $x+y+z+t=0$ et $H_2$ : $x-y+z-t=0$, et $u=(1,1,1,1)$.
\begin{enumerate}
\item Donner la dimension et une base de $H_1\cap H_2$.
\item La droite $\vect(u)$ est-elle un supplémentaire de $H_1$ ? de $H_2$ ?
\end{enumerate}
\place
\exercice{5}{1}{Projecteurs associés}{Exercice complémentaire}{projecteur, somme directe}{aln-12}{2}
Soient $F$ et $G$ deux sous-espaces supplémentaires de $E$, $p$ la projection sur $F$ parallèlement à $G$ et $q$ la projection sur
$G$ parallèlement à $F$. Montrer que $p+q=\Id_E$, $p\circ q=q\circ p=0$, $\Ker p=\Im q$ et $\Ker q=\Im p$.
\place
\exercice{6}{1}{Fonctions paires et impaires}{Exercice complémentaire}{somme directe, dimension infinie, symétrie}{rdn-17}{2}
Soit $E$ l'espace des fonctions de $\R$ dans $\R$, $\mathcal P$ le sous-espace des fonctions paires et $\mathcal I$ celui des
fonctions impaires. Montrer que $E=\mathcal P\oplus\mathcal I$. Que donne cette décomposition pour $\exp$ ?
\place
\clearpage
\partie{II}{Exercices classiques et méthodes incontournables}{Projection et symétrie en coordonnées, trace, rang d'une matrice à paramètres, projecteur donné par sa matrice, blocs.}
\exercice{7}{2}{Projecteur et symétrie dans $\mathbb R^3$}{Livret, exercice 1.2}{somme directe, projecteur, symétrie, hyperplan}{al-proj}{3}
Dans $\R^3$, on considère le plan $P$ d'équation $2x-y+z=0$ et la droite $D$ d'équations $y=0$, $x=z$.
\begin{enumerate}
\item Donner une base de $P$ et de $D$, puis montrer que $\R^3=P\oplus D$.
\item Déterminer la matrice, dans la base canonique, de la projection $p$ sur $P$ parallèlement à $D$, puis celle de la symétrie $s$
par rapport à $P$ parallèlement à $D$. Vérifier que $p\circ p=p$ et $s\circ s=\Id$.
\end{enumerate}
\place
\exercice{8}{2}{Deux propriétés de la trace}{Livret, exercice 1.4}{trace, $\mathrm{tr}(AB)=\mathrm{tr}(BA)$, projecteur}{red-trace}{3}
\begin{enumerate}
\item Montrer qu'il n'existe pas de matrices $A,B\in\Mn{\K}$ telles que $AB-BA=I_n$.
\item Soit $A\in\Mn{\K}$ telle que $A^2=A$. Montrer que $\rg(A)=\tr(A)$.
\end{enumerate}
\place
\exercice{9}{2}{La trace détermine une matrice}{Exercice complémentaire}{trace, base canonique de $\mathcal M_n$}{aln-5}{3}
Soient $A,B\in\Mn{\K}$. Montrer que $\tr(AM)=\tr(BM)$ pour tout $M\in\Mn{\K}$ si et seulement si $A=B$.
\place
\exercice{10}{2}{Rang d'une matrice à paramètres}{Exercice complémentaire}{rang, opérations élémentaires, discussion}{aln-6}{4}
Déterminer, selon les valeurs des réels $a$ et $b$, le rang de $M=\begin{pmatrix}a&b&1\\1&ab&1\\1&b&a\end{pmatrix}$.
\place
\exercice{11}{2}{Un projecteur donné par sa matrice}{Exercice complémentaire}{projecteur, rang et trace, équation d'un plan, base adaptée}{aln-4}{4}
Soit $u$ l'endomorphisme de $\R^3$ de matrice $A=\begin{pmatrix}1/2&1&-1/2\\0&1&0\\-1/2&1&1/2\end{pmatrix}$ dans la base canonique.
\begin{enumerate}
\item Montrer que $u$ est un projecteur. Quel est son rang ?
\item Déterminer une base de $\Ker u$ et de $\Im u$, puis une équation cartésienne de $\Im u$.
\item Déterminer $P$ inversible telle que $P^{-1}AP=\mathrm{diag}(0,1,1)$.
\end{enumerate}
\place
\exercice{12}{2}{Lire des blocs}{Livret, exercice 1.5}{sous-espace stable, matrice par blocs, base adaptée, trace}{red-stable}{4}
Soit $f$ l'endomorphisme de $\R^3$ de matrice $A=\begin{pmatrix}3&-1&1\\0&2&0\\1&-1&3\end{pmatrix}$ dans la base canonique
$(e_1,e_2,e_3)$.
\begin{enumerate}
\item Montrer que le plan $F$ d'équation $y=0$ et la droite $G=\vect(u)$, $u=(1,1,0)$, sont stables par $f$.
\item Montrer que $\R^3=F\oplus G$ et écrire la matrice de $f$ dans la base $(e_1,e_3,u)$. Vérifier sur cette matrice que $\tr A=8$.
\end{enumerate}
\place
\clearpage
\partie{III}{Exercices de synthèse et approfondissement}{Formule de Grassmann, matrices de rang 1, noyau et image, l'endomorphisme $M\mapsto AM-MA$, matrices nilpotentes de taille 2.}
\exercice{13}{2}{Deux sous-espaces trop grands se coupent}{Exercice complémentaire}{formule de Grassmann, intersection}{aln-3}{5}
Soient $F$ et $G$ deux sous-espaces d'un espace $E$ de dimension finie tels que $\dim F+\dim G>\dim E$. Montrer que
$F\cap G\neq\{0\}$. Application : deux plans vectoriels de $\R^3$ ont toujours une droite en commun.
\place
\exercice{14}{3}{Matrices de rang au plus $1$}{Exercice complémentaire}{rang, produit colonne-ligne, trace}{aln-7}{5}
Soit $M\in\Mn{\R}$. Montrer que $\rg M\leq1$ si et seulement s'il existe deux colonnes $U,V\in\mathcal M_{n,1}(\R)$ telles que
$M=UV^T$. Montrer alors que $M^2=\tr(M)\,M$.
\place
\exercice{15}{3}{Une matrice de carré nul}{Exercice complémentaire}{noyau et image, théorème du rang, base adaptée}{aln-8}{5}
Soit $f$ l'endomorphisme de $\R^3$ de matrice $A=\begin{pmatrix}1&1&-1\\-3&-3&3\\-2&-2&2\end{pmatrix}$ dans la base canonique.
\begin{enumerate}
\item Déterminer une base de $\Ker f$ et de $\Im f$. Comparer ces deux sous-espaces ; que vaut $f\circ f$ ?
\item Déterminer une base de $\R^3$ dans laquelle la matrice de $f$ n'a qu'un coefficient non nul.
\end{enumerate}
\place
\exercice{16}{3}{L'endomorphisme $M\mapsto AM-MA$}{Livret, exercice 1.6}{endomorphisme de matrices, trace, noyau et image, sous-espace stable}{al-crochet}{6}
Soit $A=\begin{pmatrix}1&1\\1&-1\end{pmatrix}$ et $\varphi:\mathcal M_2(\R)\to\mathcal M_2(\R)$, $M\mapsto AM-MA$.
\begin{enumerate}
\item Montrer que $\varphi$ est un endomorphisme et que son image est contenue dans l'hyperplan des matrices de trace nulle.
\item Écrire la matrice de $\varphi$ dans la base canonique $(E_{1,1},E_{1,2},E_{2,1},E_{2,2})$. Déterminer $\Ker\varphi$ et $\Im\varphi$.
\item Montrer que $\Ker\varphi$ et $\Im\varphi$ sont stables par $\varphi$. Sont-ils supplémentaires ?
\end{enumerate}
\place
\exercice{17}{3}{Matrices nilpotentes de taille 2}{Exercice complémentaire}{trace, déterminant $2\times2$, hyperplan, sous-espace engendré}{aln-11}{6}
Une matrice $A\in\mathcal M_2(\R)$ est \emph{nilpotente} s'il existe $k\geq1$ tel que $A^k=0$. On note $\mathcal N$ l'ensemble de ces matrices.
\begin{enumerate}
\item Montrer que toute $M\in\mathcal M_2(\R)$ vérifie $M^2-\tr(M)\,M+\det(M)\,I_2=0$.
\item Montrer que $A\in\mathcal N$ si et seulement si $\tr A=0$ et $\det A=0$.
\item Montrer que $\vect(\mathcal N)$ est l'hyperplan des matrices de trace nulle.
\end{enumerate}
\place
\exercice{18}{3}{$AB-BA=I$ en dimension infinie}{Exercice complémentaire}{trace, dimension finie, polynômes}{aln-13}{7}
Sur $E=\R[X]$, on considère $D:P\mapsto P'$ et $M:P\mapsto XP$.
\begin{enumerate}
\item Montrer que $D\circ M-M\circ D=\Id_E$.
\item Pourquoi cela ne contredit-il pas le fait qu'il n'existe pas de matrices $A,B\in\Mn{\R}$ telles que $AB-BA=I_n$ ?
\end{enumerate}
\place
\clearpage
\partie{IV}{Exercices difficiles et type concours}{Endomorphismes de trace nulle en dimension 2 : symétries et endomorphismes de carré nul.}
\exercice{19}{4}{Endomorphismes de $\mathbb C^2$ de trace nulle}{Exercice complémentaire}{symétrie, trace, théorème du rang, base adaptée}{aln-10}{7}
Soit $f$ un endomorphisme de $\C^2$ tel que $\tr f=0$.
\begin{enumerate}
\item Montrer que $f\circ f$ est une homothétie ; on note $\lambda$ son rapport.
\item On suppose $\lambda\neq0$. Montrer qu'il existe $\alpha\in\C$ et une symétrie $s$ de $\C^2$ telles que $f=\alpha s$, puis une base de
$\C^2$ dans laquelle la matrice de $f$ est $\begin{pmatrix}0&\alpha\\\alpha&0\end{pmatrix}$.
\item On suppose $\lambda=0$ et $f\neq0$. Montrer qu'il existe une base de $\C^2$ dans laquelle la matrice de $f$ est $\begin{pmatrix}0&1\\0&0\end{pmatrix}$.
\end{enumerate}
\place
\end{document}
